% ==========================
% APÊNDICE A
% ==========================
\chapter{Artefatos de Gestão do Projeto}\label{apx:artefatos}

Este apêndice apresenta os principais artefatos de gestão utilizados ao longo do projeto, com o objetivo de evidenciar a aplicação prática dos princípios do PMBOK 7ª edição em conjunto com métodos ágeis, em especial o Scrum. Os documentos aqui descritos sustentam a rastreabilidade das decisões, o planejamento incremental e o controle adaptativo do projeto.

\section{Visão do Produto}

O produto desenvolvido consiste em um sistema de gestão de casos sociais, com foco em registro estruturado de atendimentos, categorização por tags, georreferenciamento e histórico auditável de ações. O sistema busca apoiar a tomada de decisão institucional, a integração futura com ferramentas analíticas e a geração de evidências para políticas públicas.

\section{Papéis e Responsabilidades}

Os papéis do Scrum foram definidos da seguinte forma:
\begin{itemize}
    \item \textbf{Product Owner}: responsável pela priorização do backlog, definição de valor e interlocução com stakeholders institucionais.
    \item \textbf{Scrum Master}: responsável por facilitar o processo, remover impedimentos e garantir a aderência às práticas ágeis.
    \item \textbf{Time de Desenvolvimento}: responsável pela implementação técnica, testes e validação das entregas.
\end{itemize}

Todos os integrantes do grupo participaram ativamente das atividades de planejamento, execução e revisão, refletindo o caráter colaborativo exigido pelo projeto.

\section{Backlog Consolidado}

O backlog do produto foi estruturado em épicos e histórias de usuário, priorizadas com base em valor, risco e dependências técnicas. A cada sprint, o backlog foi refinado e ajustado conforme feedback obtido nas revisões.

\section{Definição de Pronto (Definition of Done)}

Uma história foi considerada pronta quando atendia simultaneamente aos seguintes critérios:
\begin{itemize}
    \item Funcionalidade implementada conforme critérios de aceitação;
    \item Dados persistidos de forma segura e consistente;
    \item Testes básicos realizados;
    \item Registro da entrega no repositório e documentação mínima atualizada.
\end{itemize}

\section{Métricas de Acompanhamento}

As métricas utilizadas incluíram velocidade do time, taxa de conclusão de histórias por sprint, lead time e acompanhamento qualitativo de riscos. Essas métricas serviram como instrumentos de aprendizado, e não como mecanismos punitivos.

---

% ==========================
% APÊNDICE B
% ==========================
\chapter{Visão Técnica e Estrutural do Sistema}\label{apx:tecnico}

Este apêndice descreve a visão técnica do sistema proposto, com foco na arquitetura conceitual, nos principais componentes e nas decisões de projeto adotadas para sustentar escalabilidade, rastreabilidade e segurança da informação.

\section{Arquitetura Conceitual}

O sistema foi concebido segundo uma arquitetura em camadas, composta por:
\begin{itemize}
    \item Camada de Apresentação: interface web para registro e consulta de casos;
    \item Camada de Aplicação: regras de negócio, controle de fluxo e validações;
    \item Camada de Dados: persistência estruturada de casos, tags, usuários e geolocalização.
\end{itemize}

Essa separação visa facilitar manutenção, evolução incremental e possíveis integrações futuras.

\section{Modelo de Dados Simplificado}

O modelo de dados contempla as seguintes entidades principais:
\begin{itemize}
    \item Caso Social;
    \item Usuário;
    \item Tag;
    \item Registro de Atendimento;
    \item Localização Geográfica.
\end{itemize}

Cada caso social mantém histórico completo de alterações, permitindo auditoria e análise temporal dos atendimentos.

\section{Georreferenciamento}

O sistema prevê o armazenamento de coordenadas geográficas associadas aos casos, possibilitando visualização em mapas e análise espacial. Essa funcionalidade amplia a capacidade analítica do sistema, permitindo identificar padrões regionais e áreas de maior demanda.

\section{Segurança e Ética no Uso de Dados}

Considerando a natureza sensível das informações tratadas, o projeto incorpora princípios de controle de acesso, rastreabilidade de ações e minimização de dados. Essas decisões estão alinhadas tanto às boas práticas de engenharia de software quanto aos princípios éticos defendidos pelo PMBOK.

\section{Evolução Futura}

Como possibilidades de evolução do sistema, destacam-se:
\begin{itemize}
    \item Integração com ferramentas de análise de dados e visualização;
    \item Uso de modelos de aprendizado de máquina para apoio à priorização de casos;
    \item Integração com sistemas institucionais existentes.
\end{itemize}

Essas extensões reforçam o caráter incremental e adaptativo do projeto.

% ==========================
% APÊNDICE C
% ==========================
\chapter{Riscos, Decisões Críticas e Trade-offs do Projeto}\label{apx:riscos}

Este apêndice registra os principais riscos identificados, as decisões difíceis enfrentadas ao longo do projeto e os trade-offs assumidos conscientemente pelo grupo. O objetivo não é listar problemas hipotéticos, mas evidenciar maturidade na gestão de incertezas e na tomada de decisão em ambientes complexos, conforme preconizado pelo PMBOK 7ª edição.

\section{Principais Riscos Identificados}

\subsection{Risco de Escopo Difuso}

A natureza do sistema de gestão de casos sociais envolve múltiplos perfis de usuários e necessidades institucionais distintas. Houve o risco recorrente de expansão contínua do escopo, especialmente relacionado a funcionalidades analíticas e relatórios avançados.

\textbf{Resposta adotada}: priorização explícita de um núcleo funcional mínimo, com registro estruturado de casos, tags e georreferenciamento, deixando funcionalidades avançadas como evolução futura do produto.

\subsection{Risco de Qualidade dos Dados}

A utilidade do sistema depende fortemente da qualidade dos dados inseridos pelos usuários. Erros de preenchimento, registros incompletos ou inconsistentes poderiam comprometer análises posteriores.

\textbf{Resposta adotada}: definição de campos obrigatórios, uso de validações básicas e incentivo a registros textuais objetivos, reconhecendo que controles excessivos poderiam reduzir a adoção do sistema.

\subsection{Risco de Segurança e Privacidade}

O tratamento de dados sensíveis de natureza social envolve riscos éticos, legais e institucionais, especialmente em relação a acessos indevidos ou uso inadequado das informações.

\textbf{Resposta adotada}: adoção de princípios de minimização de dados, controle de acesso por perfil e rastreabilidade de ações, mesmo em um cenário de protótipo acadêmico.

\subsection{Risco de Adoção pelos Usuários}

Sistemas institucionais frequentemente falham não por problemas técnicos, mas por resistência dos usuários à mudança de processos.

\textbf{Resposta adotada}: foco em uma interface simples, próxima à lógica de registro manual tradicional, reduzindo fricção cognitiva e facilitando a transição do processo físico para o digital.

\section{Decisões Difíceis Enfrentadas}

\subsection{Simplicidade versus Funcionalidade}

Uma decisão central do projeto foi limitar deliberadamente o número de funcionalidades iniciais, mesmo reconhecendo a demanda por recursos mais sofisticados.

Essa escolha privilegiou estabilidade, clareza de uso e previsibilidade de entrega, em detrimento de uma solução mais completa porém menos viável dentro do horizonte do projeto.

\subsection{Flexibilidade versus Padronização}

Foi necessário equilibrar a liberdade de registro textual com a necessidade de padronização mínima para análise futura.

Optou-se por um modelo híbrido, combinando campos estruturados com espaço para observações livres, aceitando a perda parcial de uniformidade em favor de maior expressividade do usuário.

\subsection{Planejamento versus Adaptação}

Embora houvesse a possibilidade de definir um planejamento mais detalhado desde o início, o grupo optou por manter o planejamento em nível incremental, ajustado a cada sprint.

Essa decisão reduziu a previsibilidade de longo prazo, mas aumentou a capacidade de resposta a descobertas feitas durante o desenvolvimento.

\section{Trade-offs Assumidos}

\subsection{Velocidade de Entrega versus Robustez Técnica}

Optou-se por soluções técnicas consolidadas e de menor complexidade, mesmo que menos otimizadas, priorizando entregas funcionais dentro do prazo do projeto.

Esse trade-off foi considerado aceitável diante do caráter acadêmico e exploratório da solução.

\subsection{Escalabilidade Futura versus Simplicidade Atual}

O projeto foi concebido com uma arquitetura que permite evolução futura, mas sem implementar mecanismos avançados de escalabilidade desde o início.

A decisão evitou sobreengenharia, ao custo de exigir refatorações em um cenário de crescimento real.

\subsection{Controle Centralizado versus Autonomia do Time}

A governança do projeto privilegiou autonomia do time de desenvolvimento, com decisões tomadas de forma colaborativa, em detrimento de um controle central rígido.

Esse trade-off aumentou o engajamento e a responsabilidade coletiva, alinhando-se aos princípios ágeis e aos valores do PMBOK 7ª edição.

\section{Síntese dos Aprendizados}

A análise dos riscos, decisões e trade-offs reforça que não existem soluções perfeitas em projetos reais, apenas escolhas contextualizadas. O sucesso do projeto esteve menos associado à eliminação de incertezas e mais à capacidade do grupo de reconhecê-las, discuti-las abertamente e responder de forma consciente e adaptativa.

Essa postura evidencia uma compreensão madura da gestão de projetos contemporânea, na qual frameworks servem como guias, e não como substitutos do julgamento profissional.
